% Imporditud: tools/import_eksperimendid.py
% Allikas: Eesti füüsikaolümpiaadi eksperimendi ülesannete
%   kogu (Rael Kalda, 2023), ülesanne Ü115, lahendus L115.
\setAuthor{Mihkel Kree}
\setRound{lõppvoor}
\setYear{2016}
\setNumber{G E1}
\setDifficulty{3}
\setTopic{elektriahelad}
\setVahendid{alalisvoolumootor tiivikuga; muudetava pingega vooluallikas, mis näitab tablool nii pinget kui voolutugevust; ühendusjuhtmed; laser-tahhomeeter; graafikupaberid lineaarse teljestikuga ning logaritmiliselt jaotatud teljestikuga}
\setPunktid{10}
\setAllikas{Eksperimendi kogu Ü115, lahendus lk 89}
\setLahendusseis{kasitsi}

\prob{Tiivik}
Kujutage graafiliselt tiiviku pöörlemiseks vajaliku elektrivõimsuse P (vahemikus 0,1 W kuni 2 W) sõltuvus tiiviku pöörlemissagedusest f . Teades, et sõltuvus avaldub kujul P = kf n, määrake astendaja n väärtus.

(tablool C.V.). Mootori käivitamiseks ei pruugi piisata pinge aeglasest suurendamisest, sel juhul ühendage voolujuhe hetkeks lahti ja siis kohe tagasi. Tahhomeetriga mõõtmisel hoidke all test-nuppu ning suunake laserkiir tiivikule, hoides tahhomeetrit sellest umbes 5 cm kaugusel. Tahhomeeter näitab lasertäpilt hajunud valguse heleduse võnkeperioode ühes minutis.
\textit{Kasutusjuhend:} Vooluallikal on pinge muutmiseks kaks nuppu: peen- (fine) ja jämereguleerimine (coarse). Jälgige, et vooluallikas töötaks konstantse pinge režiimis
\textit{Hoiatused:} Jälgige, et elektrivõimsus mootoris ei ületaks 2 W. Mootor võib läbi põleda! Ärge suunake laserikiirt endale ega teistele silma!

\hint

\solu
\textit{Lahendus on kogumikus lk 89. Siia ei ole seda veel ümber kirjutatud, sest originaalis on tabel või joonis, mis PDF-i tekstikihist loetavalt välja ei tule.}
\probend
